\documentclass[11pt,a4paper]{article}
\usepackage[margin=1in]{geometry}
\usepackage{amsmath,amssymb,amsthm}
\usepackage{booktabs}
\usepackage{longtable}
\usepackage{enumitem}
\usepackage{xcolor}
\usepackage{hyperref}
\usepackage{titlesec}

\definecolor{gradeAplus}{HTML}{006400}
\definecolor{gradeA}{HTML}{228B22}
\definecolor{gradeAminus}{HTML}{6B8E23}
\definecolor{gradeBplus}{HTML}{DAA520}
\definecolor{gradeB}{HTML}{CD853F}
\definecolor{hardproblem}{HTML}{8B0000}
\definecolor{quickwin}{HTML}{4169E1}

\titleformat{\section}{\Large\bfseries}{\thesection}{1em}{}
\titleformat{\subsection}{\large\bfseries}{\thesubsection}{1em}{}

\title{\textbf{Density Field Dynamics: Plan to A+ in Every Sector}\\[0.5em]
\large Synthesized from Multi-Agent Evaluation Feedback\\
\large Honest Assessment with Actionable Roadmap}
\author{DFD Research Programme}
\date{March 2026}

\begin{document}
\maketitle
\tableofcontents
\newpage

%% ============================================================
\section{Executive Summary}
%% ============================================================

Density Field Dynamics currently holds grades ranging from A+ (Parameter Economy) to B (Cosmology, Quantum Sector).
The theory's distinctive strength is \emph{breadth combined with specificity}: it derives actual numerical values from a single scalar-field action with no free parameters, covering territory from $\alpha$ to galactic rotation curves to fermion masses.
Its principal weaknesses are computational closure (no production-grade Boltzmann pipeline), formal mathematical industrialization (global existence proofs remain open), and the absence of peer-reviewed publication or confirmed experimental tests.

This document provides, for each of twelve evaluated sectors:
\begin{enumerate}
    \item An honest assessment of the current grade and why it sits where it does.
    \item The concrete standard that ``A+'' represents.
    \item Specific action items, dependencies, estimated effort, and timeline.
\end{enumerate}

\noindent Three overarching themes emerge:
\begin{itemize}
    \item \textbf{The Boltzmann pipeline is the single highest-leverage investment.} It unlocks A+ in Cosmology, strengthens Testability, and provides computational maturity. Until \texttt{mochi\_class} or \texttt{SymBoltz} exists, cosmological claims remain theoretical.
    \item \textbf{Global existence is a hard mathematics problem.} It gates both Math Foundations and Gravitational/Classical to A+. It may require collaboration with PDE specialists and could take years.
    \item \textbf{Peer review and experiment are externally gated.} No amount of internal work can substitute for published, refereed papers and laboratory confirmation.
\end{itemize}

\bigskip
\noindent\textbf{Current scorecard:}

\begin{center}
\begin{tabular}{lcc}
\toprule
\textbf{Sector} & \textbf{Current} & \textbf{Gap to A+} \\
\midrule
Parameter Economy & \textcolor{gradeAplus}{A+} & --- \\
Strong CP & \textcolor{gradeA}{A} & Small \\
Testability / Falsifiability & \textcolor{gradeA}{A} & Medium (external) \\
Gravitational / Classical & \textcolor{gradeAminus}{A$-$} & Medium--Hard \\
Galactic Dynamics / MOND & \textcolor{gradeAminus}{A$-$} & Medium \\
Microsector / $\alpha$ & \textcolor{gradeAminus}{A$-$} & Medium \\
Math Foundations & \textcolor{gradeAminus}{A$-$} & Hard \\
Fermion Masses & \textcolor{gradeBplus}{B+} & Medium--Hard \\
Computational Maturity & \textcolor{gradeBplus}{B+} & Medium (engineering) \\
Unification Reach & \textcolor{gradeBplus}{B+} & Hard \\
Cosmology & \textcolor{gradeB}{B} & Hard (pipeline) \\
Quantum Sector & \textcolor{gradeB}{B} & Medium \\
\bottomrule
\end{tabular}
\end{center}


%% ============================================================
\newpage
\section{Sector-by-Sector Plan}
%% ============================================================

%% ------------------------------------------------------------
\subsection{Parameter Economy --- Current Grade: A+}
%% ------------------------------------------------------------

\paragraph{Honest assessment.}
This is DFD's crown jewel. The theory has \emph{zero} freely adjustable parameters: everything---$\alpha$, fermion masses, $\kappa$, cosmological observables---derives from a single scalar-field action $S[\psi]$ plus the integers $(n, k, \ell)$ that label topological sectors. No competing framework achieves this. The Standard Model has $\sim$19 free parameters; $\Lambda$CDM adds 6 more. String Theory has a landscape of $10^{500}$ vacua with no selection principle.

\paragraph{What A+ means.}
Every derived quantity traces to the action with no tuning. Already satisfied.

\paragraph{Action items.}
\begin{enumerate}
    \item \textbf{Maintain:} As new sectors are derived (CKM, PMNS, $\Lambda_{\text{QCD}}$), ensure no free parameters creep in. Each new derivation must close with a parameter count showing zero additions.
    \item \textbf{Document:} Create a single-page ``parameter census'' table listing every quantity DFD derives, with the equation number tracing it to $S[\psi]$. Useful for referees.
\end{enumerate}

\paragraph{Effort:} 1--2 days for the census table. \quad \textbf{Dependencies:} None. \quad \textbf{Timeline:} Immediate.

\paragraph{Risk of regression:} Low, unless a future derivation requires an \emph{ad hoc} input.

%% ------------------------------------------------------------
\subsection{Strong CP --- Current Grade: A}
%% ------------------------------------------------------------

\paragraph{Honest assessment.}
DFD resolves the Strong CP problem at theorem grade: the $\psi$-field action has no $\bar{\theta}$ term, so $\bar{\theta} = 0$ is not fine-tuned but \emph{structural}. The proof is clean, the logic is short, and no known gap exists. The only reason this is not A+ is the absence of peer review confirming the argument, and the fact that no experiment has yet tested DFD's specific prediction for the neutron EDM versus axion-based competitors.

\paragraph{What A+ means.}
A refereed publication confirming the $\bar{\theta} = 0$ theorem, plus either (a) experimental confirmation that no axion exists (consistent with DFD) or (b) a lattice-QCD-level computation of the neutron EDM from DFD matching data.

\paragraph{Action items.}
\begin{enumerate}
    \item \textbf{Write a focused Strong CP paper} (standalone, 8--12 pages) presenting the $\bar{\theta}$ theorem cleanly with all lemmas explicit. Target: Physical Review D or JHEP. \emph{Effort: 2--3 weeks.}
    \item \textbf{Sharpen the neutron EDM prediction.} Current DFD predicts $d_n = 0$ at tree level; compute the one-loop correction from $\psi$-field dressing to give a precise numerical bound. \emph{Effort: 1--2 weeks of calculation.}
    \item \textbf{Compare explicitly with axion models.} Table showing: axion window vs.\ DFD prediction for $\bar{\theta}$, neutron EDM, and new CP-violating observables. \emph{Effort: 1 week.}
    \item \textbf{Submit and respond to referee reports.} \emph{Effort: 2--6 months (externally gated).}
\end{enumerate}

\paragraph{Dependencies:} None---this is self-contained. \quad \textbf{Timeline:} Paper draft in 1 month; submission by month 2; A+ contingent on peer review (6--12 months).

\paragraph{Likelihood of reaching A+:} \textcolor{gradeA}{High}. The gap is publication, not physics.

%% ------------------------------------------------------------
\subsection{Testability / Falsifiability --- Current Grade: A}
%% ------------------------------------------------------------

\paragraph{Honest assessment.}
DFD makes numerous concrete, falsifiable predictions: $\alpha$ to 9+ digits, fermion mass ratios, neutrino mass sum $\Sigma m_\nu$, $S_8$ scale-dependent running, tensor-to-scalar ratio $r$, galactic rotation curves without dark matter, and specific CMB signatures. The theory is \emph{highly} falsifiable---more so than most competitors. The gap to A+ is that \textbf{none of these predictions have been tested against DFD-specific outputs by an independent group}, and no peer-reviewed publication exists establishing the predictions formally.

The three remaining YELLOW predictions ($\Sigma m_\nu$, $S_8$ scale dependence, $r$) are all experiment-limited: the data to confirm or refute them either does not exist yet or has not been analyzed in the DFD framework.

\paragraph{What A+ means.}
At least one DFD-specific prediction confirmed by independent experiment or observation, plus peer-reviewed publication of the prediction \emph{prior} to confirmation (to rule out post-diction).

\paragraph{Action items.}
\begin{enumerate}
    \item \textbf{Publish predictions with timestamps.} Zenodo deposits with DOIs for: $\Sigma m_\nu$ prediction, $S_8(k)$ running curve, $r$ prediction, rotation curve fits. This creates a public, dated record. \emph{Effort: 1--2 weeks.}
    \item \textbf{Submit to peer review.} At minimum, the $\alpha$ derivation paper and the cosmological predictions paper. Refereed publication is the gold standard for establishing predictions. \emph{Effort: 2--6 months (externally gated).}
    \item \textbf{Engage experimentalists.} Identify groups working on:
    \begin{itemize}
        \item KATRIN / Project 8 (neutrino mass) --- share $\Sigma m_\nu$ prediction.
        \item DES / Euclid / Rubin LSST ($S_8$) --- share scale-dependent $S_8$ curve.
        \item CMB-S4 / LiteBIRD ($r$) --- share tensor-to-scalar prediction.
        \item SRF cavity groups (SQMS at Fermilab) --- share $\psi$-field coupling prediction.
    \end{itemize}
    \emph{Effort: Ongoing outreach, 1--2 months to initiate.}
    \item \textbf{Pre-register predictions.} Use OSF or similar platform to pre-register specific numerical predictions before experimental results are published. \emph{Effort: 1 week.}
\end{enumerate}

\paragraph{Dependencies:} Cosmology predictions require at least preliminary Boltzmann pipeline results. SRF predictions are self-contained.

\paragraph{Timeline:} Timestamps and Zenodo deposits: immediate. Peer review: 6--12 months. Experimental confirmation: 1--5 years (externally gated by experimental timelines).

\paragraph{Likelihood of reaching A+:} \textcolor{gradeBplus}{Medium}. The physics is there; the gate is external validation. This \emph{cannot} be rushed.

%% ------------------------------------------------------------
\subsection{Gravitational / Classical Sector --- Current Grade: A$-$}
%% ------------------------------------------------------------

\paragraph{Honest assessment.}
DFD reproduces Newtonian gravity and GR in appropriate limits. The $\psi$-field equation reduces to the Poisson equation in the weak-field, non-relativistic regime. Schwarzschild-like solutions exist. The gap to A+ is twofold:
\begin{enumerate}[label=(\alph*)]
    \item \textbf{Global existence} of solutions to the full nonlinear $\psi$-field PDE is unproven. Local well-posedness is established; global existence is open.
    \item \textbf{Horizon boundary conditions:} The treatment of black hole horizons in the $\psi$-field framework is incomplete. How does $\psi$ behave at $r \to r_s$? Does the field remain regular? What replaces the Penrose diagram?
\end{enumerate}

\paragraph{What A+ means.}
A complete existence and uniqueness theory for the $\psi$-field equation in physically relevant geometries, including black hole spacetimes, with horizon boundary conditions rigorously specified. Plus: post-Newtonian corrections computed to 2PN or higher, matching GR predictions for binary pulsars and gravitational wave templates.

\paragraph{Action items.}
\begin{enumerate}
    \item \textbf{Post-Newtonian expansion to 2PN.} Derive the $\psi$-field corrections to GR at 2PN order. Compare with Hulse--Taylor binary pulsar timing. This is achievable with existing tools. \emph{Effort: 2--4 weeks.}
    \item \textbf{Gravitational wave templates.} Compute inspiral waveforms in DFD and compare with LIGO/Virgo templates. Even showing consistency at leading order is valuable. \emph{Effort: 1--2 months.}
    \item \textbf{Horizon boundary conditions.} Formulate the $\psi$-field regularity condition at horizons. Does $\psi \to \psi_0$ (finite)? Does $\nabla\psi$ diverge? Classify solutions. \emph{Effort: 2--4 weeks of analysis.}
    \item \textbf{Global existence program.} This is a \emph{hard mathematics problem}. Options:
    \begin{itemize}
        \item Prove global existence for the spherically symmetric case first (tractable).
        \item Seek energy estimates or monotonicity formulae that prevent blow-up.
        \item Collaborate with PDE analysts (e.g., groups working on wave maps, Yang--Mills).
    \end{itemize}
    \emph{Effort: 6 months to years. See Hard Problems section.}
\end{enumerate}

\paragraph{Dependencies:} Items 1--3 are independent. Item 4 depends on Math Foundations progress.

\paragraph{Timeline:} Post-Newtonian and GW templates: 1--3 months. Horizon BCs: 1--2 months. Global existence: 6 months--3 years.

\paragraph{Likelihood of reaching A+:} \textcolor{gradeBplus}{Medium}. Items 1--3 are achievable and would bring the grade to A. Item 4 (global existence) is the hard gate to A+.

%% ------------------------------------------------------------
\subsection{Galactic Dynamics / MOND Sector --- Current Grade: A$-$}
%% ------------------------------------------------------------

\paragraph{Honest assessment.}
DFD reproduces MOND-like phenomenology (flat rotation curves, Tully--Fisher relation) from first principles without dark matter. The $\psi$-field's self-interaction naturally produces the $a_0$ acceleration scale. However, there is a \textbf{Bullet Cluster deficit of 2--3$\times$}: the observed gravitational lensing in the Bullet Cluster requires more mass separation than DFD's $\psi$-field alone provides. This is the single clearest gap.

\paragraph{What A+ means.}
Rotation curves fitted to the quality of SPARC data (150+ galaxies, $<$10\% residuals), Bullet Cluster lensing reproduced within observational errors, and the radial acceleration relation (RAR) derived from first principles matching McGaugh et al.\ data.

\paragraph{Action items.}
\begin{enumerate}
    \item \textbf{Bullet Cluster resolution.} Three strategies, in order of promise:
    \begin{enumerate}[label=(\roman*)]
        \item \textbf{Non-equilibrium $\psi$-field dynamics.} The Bullet Cluster is a collision---solve the time-dependent $\psi$-field equation for merging clusters. The offset may emerge from transient dynamics. \emph{Effort: 2--4 weeks of numerical work.}
        \item \textbf{Higher-order $\psi$-field coupling.} Check whether $\psi^4$ or $\psi^6$ terms in the action provide additional lensing convergence at cluster scales. \emph{Effort: 1--2 weeks.}
        \item \textbf{Baryonic inventory.} Re-examine whether the hot gas mass in the Bullet Cluster has been underestimated. Some MOND analyses suggest the deficit may be smaller than 2--3$\times$. \emph{Effort: 1 week literature review.}
    \end{enumerate}
    \item \textbf{SPARC galaxy fits.} Run DFD rotation curve predictions against the full SPARC sample (175 galaxies). Publish residuals. \emph{Effort: 2--3 weeks with existing code.}
    \item \textbf{RAR derivation.} Show analytically that the DFD $\psi$-field equation implies $g_{\text{obs}} = g_{\text{bar}} / (1 - e^{-\sqrt{g_{\text{bar}}/g_\dagger}})$ with $g_\dagger$ derived from the action. \emph{Effort: 1--2 weeks.}
    \item \textbf{Cluster-scale predictions.} Beyond Bullet Cluster: predict lensing profiles for Abell 520 (``train wreck'' cluster), El Gordo, and other merging clusters. \emph{Effort: 2--3 weeks.}
\end{enumerate}

\paragraph{Dependencies:} Items 2--4 are independent. Item 1(i) benefits from improved numerical solvers (Computational Maturity).

\paragraph{Timeline:} SPARC fits and RAR: 1--2 months. Bullet Cluster: 1--3 months. Cluster predictions: 2--3 months.

\paragraph{Likelihood of reaching A+:} \textcolor{gradeBplus}{Medium}. The Bullet Cluster is a genuine challenge for any MOND-like theory. If the non-equilibrium dynamics solve it, A+ is achievable. If not, this sector may plateau at A.

%% ------------------------------------------------------------
\subsection{Microsector / $\alpha$ --- Current Grade: A$-$}
%% ------------------------------------------------------------

\paragraph{Honest assessment.}
The derivation of $\alpha \approx 1/137.036$ from the $\psi$-field action is DFD's most striking result. The derivation proceeds through a chain of $\sim$12 steps, and independent numerical verification confirms the output to 9+ significant figures. However, two steps have drawn scrutiny:
\begin{enumerate}[label=(\alph*)]
    \item \textbf{Step 9: the $B/A$ ratio.} The ratio of two spectral coefficients is asserted to equal a specific value. The derivation relies on a Selberg-type trace formula identity that, while numerically verified, lacks a fully self-contained proof from the action alone.
    \item \textbf{$k_{\max}$ bundle origin.} The truncation of the topological sum at $k_{\max} = 60$ (or similar) is motivated by physical arguments (stability of the $\psi$-field mode) but the mathematical proof that higher modes are absent is incomplete.
\end{enumerate}

\paragraph{What A+ means.}
Every step in the $\alpha$ derivation chain proven from the action with no gaps, $k_{\max}$ derived (not assumed), and independent reproduction by at least one external group.

\paragraph{Action items.}
\begin{enumerate}
    \item \textbf{Close Step 9.} Write a self-contained proof of the $B/A$ ratio from the spectral geometry of the $\psi$-field on $S^3 \times \mathbb{R}$. Two approaches:
    \begin{enumerate}[label=(\roman*)]
        \item Direct computation via heat kernel coefficients $a_0, a_2, a_4$ on $S^3$.
        \item Selberg trace formula applied to the relevant Laplacian, with explicit error bounds.
    \end{enumerate}
    \emph{Effort: 2--4 weeks of focused mathematical work.}
    \item \textbf{Derive $k_{\max}$ from stability.} Show that modes with $k > k_{\max}$ are dynamically unstable (exponentially growing) or violate the energy positivity condition of the action. This should follow from the second variation $\delta^2 S[\psi]$. \emph{Effort: 1--2 weeks.}
    \item \textbf{Package the derivation for external verification.} Create a Mathematica/Python notebook that takes the action as input and outputs $\alpha$ with every intermediate step annotated. This enables any physicist to verify the chain independently. \emph{Effort: 1--2 weeks (partly done---the Julia code exists).}
    \item \textbf{Submit the $\alpha$ paper to peer review.} The current PDF is detailed but unpublished. Target: Physical Review Letters (letter) + Physical Review D (companion). \emph{Effort: 2--4 weeks for reformatting; 6--12 months for review.}
\end{enumerate}

\paragraph{Dependencies:} Step 9 and $k_{\max}$ are mathematically independent. External verification requires the notebook (item 3).

\paragraph{Timeline:} Step 9 proof and $k_{\max}$ derivation: 1--2 months. Verification notebook: 1 month. Peer review: 6--12 months.

\paragraph{Likelihood of reaching A+:} \textcolor{gradeA}{High}. The gaps are well-defined and the numerical verification gives confidence that the proofs exist. This is one of the most achievable upgrades.

%% ------------------------------------------------------------
\subsection{Math Foundations --- Current Grade: A$-$}
%% ------------------------------------------------------------

\paragraph{Honest assessment.}
DFD's mathematical framework is well-defined: the action is specified, the field equation is a nonlinear elliptic/parabolic PDE, variational structure is established, and local well-posedness follows from standard theory. Spectral geometry on $S^3$ is used rigorously. The gaps are:
\begin{enumerate}[label=(\alph*)]
    \item \textbf{Global existence and uniqueness.} The full nonlinear $\psi$-field equation has not been shown to have global solutions in 3+1 dimensions. This is a hard open problem, comparable in difficulty to global regularity for Navier--Stokes or Yang--Mills (both Millennium Prize problems).
    \item \textbf{Horizon boundary conditions.} Same as Gravitational sector.
    \item \textbf{Leray--Schauder branch analysis.} The topological branching structure (which branches correspond to physical vacua) is partially analyzed but not complete.
\end{enumerate}

\paragraph{What A+ means.}
Global existence and uniqueness for the $\psi$-field equation (at least for small data or symmetric configurations), rigorous horizon boundary conditions, and a complete branch classification.

\paragraph{Action items.}
\begin{enumerate}
    \item \textbf{Spherically symmetric global existence.} Reduce the PDE to 1+1 dimensions under spherical symmetry and prove global existence using energy methods. This is the most tractable sub-problem. \emph{Effort: 1--3 months.}
    \item \textbf{Small-data global existence in 3+1.} For data close to the ground state $\psi_0$, use perturbative techniques (bootstrap/continuity argument) to prove global existence. \emph{Effort: 3--6 months.}
    \item \textbf{Energy monotonicity.} Find a Lyapunov functional or monotonicity formula that prevents finite-time blow-up. If one exists, global existence follows. If not, document what goes wrong. \emph{Effort: 2--4 months.}
    \item \textbf{Branch classification.} Complete the Leray--Schauder analysis: enumerate all branches of the nonlinear eigenvalue problem, prove which are stable, and show the physical vacuum is the unique stable branch. \emph{Effort: 2--3 months.}
    \item \textbf{External collaboration.} Engage mathematicians working on nonlinear PDE (geometric analysis, dispersive PDE). The $\psi$-field equation may be amenable to techniques from wave maps or harmonic map heat flow. \emph{Effort: Ongoing.}
\end{enumerate}

\paragraph{Dependencies:} Item 1 is a prerequisite for item 2. Item 3 is independent and could shortcut items 1--2 if successful.

\paragraph{Timeline:} Spherically symmetric case: 1--3 months. Small data: 3--6 months. Full global existence: 1--5 years (possibly open-ended).

\paragraph{Likelihood of reaching A+:} \textcolor{hardproblem}{Low--Medium}. This is one of the hardest gaps. Spherically symmetric global existence is likely achievable and would bring the grade to A. Full 3+1 global existence may require a breakthrough or may remain open indefinitely---as it has for Yang--Mills and Navier--Stokes.

%% ------------------------------------------------------------
\subsection{Quantum Sector --- Current Grade: B}
%% ------------------------------------------------------------

\paragraph{Honest assessment.}
DFD claims to resolve quantum measurement via the $\psi$-field's intrinsic decoherence mechanism: the scalar field provides an objective collapse/branching trigger without needing an external observer. The Page--Geilker experiment (which tested gravitational effects of quantum superposition) has been analyzed in agent work and the resolution \emph{exists}, but it is \textbf{not formalized in any published or even drafted paper}. The branching picture (how and when superpositions resolve into definite outcomes) is described qualitatively but lacks the mathematical precision of, say, decoherent histories or GRW theory.

\paragraph{What A+ means.}
A complete, mathematically precise measurement theory: state reduction axioms derived from the action, decoherence rates computed and compared with experiment, Born rule derived (not assumed), Page--Geilker resolved in a published paper, and the theory's position relative to Everett/Copenhagen/GRW/objective collapse clearly delineated.

\paragraph{Action items.}
\begin{enumerate}
    \item \textbf{Formalize the Page--Geilker resolution.} The analysis exists in agent research files. Turn it into a self-contained 10--15 page paper. Key content: state the DFD decoherence mechanism, apply it to the Page--Geilker setup, compute the predicted outcome, compare with observation. \emph{Effort: 2--3 weeks.}
    \item \textbf{Derive decoherence rates.} From the $\psi$-field coupling to matter, compute:
    \begin{itemize}
        \item Decoherence rate as a function of mass and spatial separation.
        \item Comparison with observed decoherence in mesoscopic systems (SQUIDs, fullerene interferometry).
        \item Prediction for future experiments (e.g., MAQRO satellite mission).
    \end{itemize}
    \emph{Effort: 2--4 weeks.}
    \item \textbf{Born rule derivation.} Show that $P = |\langle \psi_{\text{branch}} | \Psi \rangle|^2$ follows from the $\psi$-field dynamics, not as an axiom but as a consequence. This may require the branching structure to be fully specified. \emph{Effort: 1--3 months (hard).}
    \item \textbf{Position paper: DFD vs.\ interpretations.} Write a clear comparison showing how DFD relates to:
    \begin{itemize}
        \item Copenhagen (DFD provides the ``observer'' via $\psi$-field decoherence).
        \item Many-Worlds (DFD branches are selected by $\psi$-field dynamics, not all realized).
        \item GRW/CSL (DFD provides the collapse mechanism from the action, not \emph{ad hoc}).
        \item Penrose objective reduction (closest analog---but DFD derives the threshold).
    \end{itemize}
    \emph{Effort: 2--3 weeks.}
\end{enumerate}

\paragraph{Dependencies:} Item 1 is self-contained and should be done first. Items 2--4 can proceed in parallel. Item 3 (Born rule) is the hardest and may depend on progress in Math Foundations (branch classification).

\paragraph{Timeline:} Page--Geilker paper: 1 month. Decoherence rates: 1--2 months. Born rule: 3--12 months. Position paper: 1 month.

\paragraph{Likelihood of reaching A+:} \textcolor{gradeBplus}{Medium}. Formalizing the Page--Geilker resolution and computing decoherence rates would bring this to B+/A$-$. The Born rule derivation is the hard gate. Even some well-established interpretations (Many-Worlds, for instance) have not derived the Born rule satisfactorily, so achieving this would be exceptional.

%% ------------------------------------------------------------
\subsection{Fermion Masses --- Current Grade: B+}
%% ------------------------------------------------------------

\paragraph{Honest assessment.}
DFD derives fermion mass ratios from topological quantum numbers $(n, k, \ell)$ in the $\psi$-field spectrum. This is remarkable: no other framework derives quark and lepton masses from first principles. However:
\begin{enumerate}[label=(\alph*)]
    \item \textbf{The $n_f$ exponents are at Grade C+.} The family-dependent exponents that distinguish generations (e.g., electron vs.\ muon vs.\ tau) are motivated but not rigorously derived from the action. They currently look more like a successful fit than a derivation.
    \item \textbf{The electron mass has a 3.3\% error.} While impressive for a first-principles calculation, this is an order of magnitude worse than the precision achieved for $\alpha$. The source of the error is unclear: is it a truncation effect, a missing radiative correction, or a structural issue?
    \item \textbf{QCD running.} Mass predictions at the $\psi$-field scale must be run to physical scales using QCD beta functions. The running is standard, but the matching conditions at the $\psi$-field scale are not fully specified.
\end{enumerate}

\paragraph{What A+ means.}
All 12 fermion masses (6 quarks + 3 charged leptons + 3 neutrinos) derived to $<$1\% accuracy with every exponent and quantum number assignment proven from the action.

\paragraph{Action items.}
\begin{enumerate}
    \item \textbf{Derive $n_f$ exponents from spectral geometry.} The generation index should emerge from the eigenvalue structure on $S^3$ or $\mathbb{CP}^2$. Specifically:
    \begin{itemize}
        \item Show that the multiplicity structure of the Dirac spectrum on the internal manifold yields exactly 3 generations.
        \item Prove that the exponents are $n_f = \{n_1, n_2, n_3\}$ from the eigenvalue ratios.
    \end{itemize}
    \emph{Effort: 1--3 months.}
    \item \textbf{Diagnose the 3.3\% electron error.} Systematically check:
    \begin{itemize}
        \item Radiative corrections (QED self-energy at the $\psi$-field scale).
        \item Higher-order heat kernel coefficients ($a_6$ and beyond).
        \item Truncation of the topological sum.
        \item Scheme-dependence of the matching condition.
    \end{itemize}
    \emph{Effort: 2--4 weeks.}
    \item \textbf{Specify matching conditions.} Define precisely at what scale the DFD prediction ``hands off'' to standard QCD/QED running, and compute the threshold corrections. \emph{Effort: 1--2 weeks.}
    \item \textbf{Neutrino masses.} Currently predicted as $\Sigma m_\nu \sim 0.06$ eV (normal ordering). Derive the individual masses $m_1, m_2, m_3$ and the ordering from the action. \emph{Effort: 2--4 weeks.}
    \item \textbf{Heavy quark masses.} The top quark mass prediction should be sharpened and compared with the latest Tevatron/LHC combination. \emph{Effort: 1--2 weeks.}
\end{enumerate}

\paragraph{Dependencies:} Item 1 ($n_f$ exponents) is the critical blocker. Items 2--5 can proceed in parallel.

\paragraph{Timeline:} Electron error diagnosis: 1 month. $n_f$ derivation: 1--3 months. Full mass table: 3--6 months.

\paragraph{Likelihood of reaching A+:} \textcolor{gradeBplus}{Medium}. If the $n_f$ exponents can be derived rigorously, this sector jumps to A or A$-$ immediately. The 3.3\% electron error is likely fixable. A+ requires everything to close cleanly.

%% ------------------------------------------------------------
\subsection{Cosmology --- Current Grade: B}
%% ------------------------------------------------------------

\paragraph{Honest assessment.}
This is DFD's biggest weakness relative to $\Lambda$CDM. The issue is not that DFD's cosmological predictions are wrong---early results from the Bolt.jl computation are promising---but that \textbf{DFD lacks a production-grade Boltzmann solver}. $\Lambda$CDM has a 20+ year head start: CLASS, CAMB, CosmoMC, Cobaya, MontePython are all mature, validated, and trusted by the community. DFD has preliminary Julia code (Bolt.jl modifications) that has computed a first CMB spectrum, but this is not production-grade. Specific gaps:
\begin{enumerate}[label=(\alph*)]
    \item \textbf{$\Sigma m_\nu$ pressure:} The neutrino mass sum prediction ($\sim$0.06 eV) is in tension with some analyses but consistent with KATRIN bounds. The DFD-specific signature in the CMB has not been computed.
    \item \textbf{$H_0$ tension:} DFD claims to resolve the Hubble tension, but the quantitative prediction has not been computed through a full Boltzmann pipeline to compare with SH0ES, Planck, and DESI simultaneously.
    \item \textbf{Boltzmann closure:} The $\psi$-field perturbation equation must be consistently coupled into the Boltzmann hierarchy. This is specified in the equations but not implemented in running code.
\end{enumerate}

\paragraph{What A+ means.}
A full, validated Boltzmann solver (\texttt{mochi\_class} or \texttt{SymBoltz}) that:
\begin{itemize}
    \item Reproduces the Planck CMB TT/TE/EE power spectra within statistical error.
    \item Predicts the BAO scale consistent with DESI observations.
    \item Resolves the $H_0$ tension quantitatively.
    \item Predicts the matter power spectrum $P(k)$ consistent with galaxy surveys.
    \item Passes all standard consistency checks (energy conservation, gauge invariance, etc.).
\end{itemize}

\paragraph{Action items.}
\begin{enumerate}
    \item \textbf{Build \texttt{mochi\_class} (modified CLASS with $\psi$-field).}
    \begin{itemize}
        \item Fork CLASS (C code) or hi\_class.
        \item Add $\psi$-field background equation to the Friedmann solver.
        \item Add $\psi$-field perturbation to the Boltzmann hierarchy.
        \item Validate against known limits ($\psi \to 0$ recovers $\Lambda$CDM).
    \end{itemize}
    \emph{Effort: 3--6 months of dedicated coding.}
    \item \textbf{Alternatively: build \texttt{SymBoltz} (symbolic Boltzmann in Julia).}
    \begin{itemize}
        \item Extend the existing Bolt.jl work to production quality.
        \item Implement all standard physics: recombination, reionization, lensing, neutrinos.
        \item Add $\psi$-field perturbations.
    \end{itemize}
    \emph{Effort: 4--8 months.}
    \item \textbf{CMB power spectrum comparison.} With either tool, compute $C_\ell^{TT}$, $C_\ell^{TE}$, $C_\ell^{EE}$ and overlay with Planck 2018 data. Publish residuals. \emph{Effort: 1--2 months after pipeline exists.}
    \item \textbf{$H_0$ computation.} Run the DFD Friedmann equation forward to $z=0$ and extract $H_0$. Compare with SH0ES ($73.0 \pm 1.0$) and Planck ($67.4 \pm 0.5$). \emph{Effort: 1--2 weeks after pipeline exists.}
    \item \textbf{Matter power spectrum.} Compute $P(k)$ and compare with SDSS/DESI. Check $S_8$ at multiple scales. \emph{Effort: 2--4 weeks after pipeline exists.}
    \item \textbf{MCMC parameter estimation.} Hook the Boltzmann solver into a sampler (Cobaya, emcee) and run against Planck+BAO+SN data. With zero free parameters, this is a goodness-of-fit test, not a fit. \emph{Effort: 2--4 weeks after pipeline exists.}
\end{enumerate}

\paragraph{Dependencies:} \textbf{Everything depends on item 1 or item 2.} The Boltzmann pipeline is the critical path. Items 3--6 are downstream.

\paragraph{Timeline:} Pipeline: 3--8 months. CMB comparison: month 4--9. Full validation: month 6--12.

\paragraph{Likelihood of reaching A+:} \textcolor{hardproblem}{Medium--Low in near term}. The engineering is achievable but labor-intensive. A+ requires not just a working pipeline but a \emph{validated} one that the community trusts. This likely takes 1--2 years of dedicated effort. However, even partial results (e.g., CMB first peak location, $H_0$ value) would lift the grade significantly.

%% ------------------------------------------------------------
\subsection{Computational Maturity --- Current Grade: B+}
%% ------------------------------------------------------------

\paragraph{Honest assessment.}
DFD has working code: the $\alpha$ computation runs in Julia, fermion masses are computed numerically, and a first CMB spectrum has been produced via Bolt.jl. However, the code is research-grade, not production-grade. There is no CI/CD pipeline, no comprehensive test suite, no documentation for external users, and no packaged release that someone could \texttt{pip install} or \texttt{Pkg.add} and immediately use.

\paragraph{What A+ means.}
A suite of validated, documented, tested, and publicly available codes covering all DFD predictions, usable by external researchers without hand-holding.

\paragraph{Action items.}
\begin{enumerate}
    \item \textbf{Package the $\alpha$ computation.} Turn the existing Julia code into a proper package with:
    \begin{itemize}
        \item Documentation (docstrings + Documenter.jl).
        \item Test suite with CI (GitHub Actions).
        \item A single entry point: \texttt{compute\_alpha()} that returns the value with all intermediate steps logged.
    \end{itemize}
    \emph{Effort: 1--2 weeks.}
    \item \textbf{Package the fermion mass computation.} Same treatment as above. \emph{Effort: 1--2 weeks.}
    \item \textbf{Build \texttt{mochi\_class} or \texttt{SymBoltz}.} (Same as Cosmology sector.) \emph{Effort: 3--8 months.}
    \item \textbf{Numerical validation suite.} For each derived quantity, create a regression test that checks the output against the known value. Any code change that shifts a prediction triggers a failure. \emph{Effort: 1--2 weeks.}
    \item \textbf{Zenodo/GitHub release.} Create a versioned release with DOI for each code package. \emph{Effort: 1--2 days per package.}
\end{enumerate}

\paragraph{Dependencies:} Items 1, 2, 4, 5 are independent of the Boltzmann pipeline. Item 3 is the big dependency.

\paragraph{Timeline:} $\alpha$ and fermion packages: 1 month. Boltzmann solver: 3--8 months. Full maturity: 6--12 months.

\paragraph{Likelihood of reaching A+:} \textcolor{gradeA}{High}. This is primarily engineering work, not physics. Given sufficient person-hours, A+ is achievable within a year.

%% ------------------------------------------------------------
\subsection{Unification Reach --- Current Grade: B+}
%% ------------------------------------------------------------

\paragraph{Honest assessment.}
DFD claims to be a unified theory deriving all of fundamental physics from a single action. The breadth is genuine: $\alpha$, strong CP, fermion masses, gravity, dark matter, dark energy, and cosmology all emerge from one framework. However, two significant sectors remain incomplete:
\begin{enumerate}[label=(\alph*)]
    \item \textbf{CKM/PMNS mixing matrices.} The quark mixing matrix (CKM) and neutrino mixing matrix (PMNS) are not yet derived from the action. Without these, DFD cannot claim to have derived all Standard Model parameters.
    \item \textbf{$\Lambda_{\text{QCD}}$ not independently derived.} The QCD scale parameter is used in the running of quark masses but is taken as an input, not derived from the $\psi$-field action.
\end{enumerate}

\paragraph{What A+ means.}
Every Standard Model parameter---all 19 of them---derived from the action. Plus: gravity, dark matter, dark energy, and neutrino masses. The CKM matrix elements, the PMNS matrix elements, $\Lambda_{\text{QCD}}$, and $\alpha_s(M_Z)$ all computed to percent-level accuracy.

\paragraph{Action items.}
\begin{enumerate}
    \item \textbf{Derive the CKM matrix.} The mixing angles should emerge from the overlap integrals of different-generation eigenfunctions on the internal manifold. Specifically:
    \begin{itemize}
        \item Compute $\langle u_i | d_j \rangle$ on the spectral manifold for $i,j = 1,2,3$.
        \item Show the resulting $3\times3$ unitary matrix matches the Wolfenstein parametrization.
        \item Predict $|V_{us}|$, $|V_{cb}|$, $|V_{ub}|$ to percent level.
    \end{itemize}
    \emph{Effort: 2--4 months (hard).}
    \item \textbf{Derive the PMNS matrix.} Same approach for the neutrino sector. The large mixing angles ($\theta_{12} \sim 34°$, $\theta_{23} \sim 45°$, $\theta_{13} \sim 8.5°$) must emerge. \emph{Effort: 2--4 months (hard, and depends on neutrino mass derivation).}
    \item \textbf{Derive $\Lambda_{\text{QCD}}$.} This requires computing $\alpha_s$ at the $\psi$-field scale from the non-Abelian sector of the action, then running it down to the QCD scale. The key is the gauge coupling unification structure. \emph{Effort: 1--3 months.}
    \item \textbf{Derive $\alpha_s(M_Z)$.} Closely related to $\Lambda_{\text{QCD}}$. If the gauge couplings unify at the $\psi$-field scale, $\alpha_s(M_Z)$ is determined by the running. \emph{Effort: 1--2 months.}
    \item \textbf{Full Standard Model parameter table.} Once all 19 parameters are derived, publish a single table comparing DFD predictions with PDG values. \emph{Effort: 1--2 weeks (after all derivations complete).}
\end{enumerate}

\paragraph{Dependencies:} CKM depends on having the quark mass derivation solid (Fermion Masses sector). PMNS depends on neutrino masses. $\Lambda_{\text{QCD}}$ depends on the gauge sector being fully derived.

\paragraph{Timeline:} $\Lambda_{\text{QCD}}$ and $\alpha_s$: 1--3 months. CKM: 2--4 months. PMNS: 3--6 months. Full table: 6--12 months.

\paragraph{Likelihood of reaching A+:} \textcolor{hardproblem}{Low--Medium}. Deriving the CKM and PMNS matrices from first principles would be an extraordinary achievement---no theory has done this. It is possible that the $\psi$-field framework does not uniquely determine the mixing angles, in which case this sector cannot reach A+ without additional structure. This is one of the sectors that \emph{may never} reach A+.

%% ============================================================
\newpage
\section{Critical Path Analysis}
%% ============================================================

The following diagram shows which upgrades unlock others. Items higher in the chain must be completed first.

\begin{verbatim}
                    [Boltzmann Pipeline]
                    /        |         \
                   /         |          \
     [CMB Spectrum]  [H0 Computation]  [P(k), S8]
          |               |                |
     [Cosmology A+]  [Testability A+]  [Comp. Maturity A+]
                          |
                   [Peer Review]
                   /           \
       [Strong CP A+]    [alpha A+]


     [Global Existence (spherical)]
              |
     [Global Existence (small data)]
              |
     [Math Foundations A+]  -->  [Gravitational A+]


     [n_f Exponent Derivation]
              |
     [Fermion Masses A]  -->  [CKM/PMNS]
                                   |
                            [Unification A+]


     [Page-Geilker Paper]
              |
     [Decoherence Rates]
              |
     [Born Rule Derivation]
              |
     [Quantum Sector A+]
\end{verbatim}

\bigskip
\noindent\textbf{The single highest-leverage item is the Boltzmann pipeline.} It directly unlocks three sectors (Cosmology, Computational Maturity, Testability) and indirectly strengthens all others by providing computational evidence.

\noindent\textbf{The second highest-leverage item is peer review.} Multiple sectors are one refereed paper away from upgrading.


%% ============================================================
\section{Quick Wins (Achievable in Weeks)}
%% ============================================================

These items require no breakthroughs and can be completed with existing knowledge and tools:

\begin{enumerate}
    \item \textcolor{quickwin}{\textbf{Parameter census table.}} One-page document listing every DFD-derived quantity with equation numbers tracing to $S[\psi]$. \emph{Effort: 1--2 days.}

    \item \textcolor{quickwin}{\textbf{Zenodo timestamp deposits.}} Upload predictions ($\Sigma m_\nu$, $S_8(k)$, $r$, rotation curves) with DOIs to establish priority. \emph{Effort: 1--2 days.}

    \item \textcolor{quickwin}{\textbf{Page--Geilker paper draft.}} The analysis exists in agent files. Formalize into a 10--15 page paper. Upgrades Quantum Sector from B to B+. \emph{Effort: 2--3 weeks.}

    \item \textcolor{quickwin}{\textbf{$\alpha$ verification notebook.}} Package the Julia computation into a clean, annotated notebook for external verification. \emph{Effort: 1--2 weeks.}

    \item \textcolor{quickwin}{\textbf{Strong CP standalone paper.}} Already near-complete in content. Reformatting for journal submission. \emph{Effort: 2--3 weeks.}

    \item \textcolor{quickwin}{\textbf{SPARC rotation curve fits.}} Run existing code against 175 SPARC galaxies and tabulate residuals. \emph{Effort: 2--3 weeks.}

    \item \textcolor{quickwin}{\textbf{Electron mass error diagnosis.}} Systematic check of radiative corrections, truncation, and matching conditions for the 3.3\% error. \emph{Effort: 2--4 weeks.}

    \item \textcolor{quickwin}{\textbf{DFD vs.\ interpretations position paper.}} Comparison of DFD's quantum measurement approach with Copenhagen, Everett, GRW, Penrose. \emph{Effort: 2--3 weeks.}

    \item \textcolor{quickwin}{\textbf{Code packaging ($\alpha$ + fermion mass).}} Add tests, documentation, and CI to existing Julia codes. \emph{Effort: 2--4 weeks.}

    \item \textcolor{quickwin}{\textbf{Pre-register predictions on OSF.}} Formal pre-registration of testable predictions before experimental results. \emph{Effort: 1 week.}
\end{enumerate}

\noindent\textbf{Total estimated time for all quick wins: 2--3 months of focused work.}
These alone would upgrade: Strong CP (A $\to$ A, closer to A+), Quantum (B $\to$ B+), Computational Maturity (B+ $\to$ A$-$), Testability (A $\to$ A, closer to A+).


%% ============================================================
\section{Hard Problems (Requiring Breakthroughs)}
%% ============================================================

These items are genuine research problems where success is not guaranteed:

\begin{enumerate}
    \item \textcolor{hardproblem}{\textbf{Global existence of $\psi$-field solutions in 3+1 dimensions.}}
    \begin{itemize}
        \item \emph{Why hard:} Comparable to Navier--Stokes or Yang--Mills regularity. The nonlinearity in the $\psi$-field equation may produce singularities.
        \item \emph{Impact:} Gates Math Foundations and Gravitational sectors to A+.
        \item \emph{Approach:} Start with symmetric reductions. Seek energy monotonicity. Engage PDE specialists.
        \item \emph{Estimated timeline:} 1--5 years, possibly open-ended.
    \end{itemize}

    \item \textcolor{hardproblem}{\textbf{Born rule derivation from $\psi$-field dynamics.}}
    \begin{itemize}
        \item \emph{Why hard:} No quantum interpretation has satisfactorily derived the Born rule from dynamics alone. Everettian attempts (Deutsch, Wallace, Zurek) remain controversial.
        \item \emph{Impact:} Gates Quantum Sector to A+.
        \item \emph{Approach:} The $\psi$-field's branching structure may provide a natural measure. Needs careful mathematical formulation.
        \item \emph{Estimated timeline:} 6 months--open-ended.
    \end{itemize}

    \item \textcolor{hardproblem}{\textbf{CKM and PMNS matrix derivation.}}
    \begin{itemize}
        \item \emph{Why hard:} No theory has ever derived quark/lepton mixing angles from first principles. The mixing matrix involves subtle interference between generation eigenstates.
        \item \emph{Impact:} Gates Unification Reach to A+.
        \item \emph{Approach:} Compute overlap integrals on the internal spectral manifold. May require the full non-commutative geometry of the $\psi$-field.
        \item \emph{Estimated timeline:} 3--12 months if the framework permits it; may be fundamentally blocked if additional structure is needed.
    \end{itemize}

    \item \textcolor{hardproblem}{\textbf{Bullet Cluster resolution.}}
    \begin{itemize}
        \item \emph{Why hard:} The 2--3$\times$ lensing deficit is a well-known problem for MOND-like theories. It may require genuinely new physics (e.g., non-equilibrium $\psi$-field dynamics creating transient lensing enhancement).
        \item \emph{Impact:} Gates Galactic Dynamics to A+.
        \item \emph{Approach:} Time-dependent numerical simulation of $\psi$-field in merging clusters.
        \item \emph{Estimated timeline:} 1--3 months for the simulation; the physics outcome is uncertain.
    \end{itemize}

    \item \textcolor{hardproblem}{\textbf{Production-grade Boltzmann pipeline.}}
    \begin{itemize}
        \item \emph{Why hard:} Not conceptually hard, but \emph{engineering}-hard. CLASS took years to develop with a dedicated team. DFD needs to modify it to include $\psi$-field perturbations without breaking the validated physics.
        \item \emph{Impact:} Gates Cosmology, Computational Maturity, and Testability sectors.
        \item \emph{Approach:} Fork CLASS/hi\_class. Incremental modifications with regression tests at each step.
        \item \emph{Estimated timeline:} 3--8 months for a working version; 1--2 years for community-trusted validation.
    \end{itemize}

    \item \textcolor{hardproblem}{\textbf{$n_f$ generation exponents from spectral geometry.}}
    \begin{itemize}
        \item \emph{Why hard:} The ``three generations'' problem is one of the deepest unsolved questions in particle physics. Deriving the generation structure from geometry is the holy grail of spectral approaches.
        \item \emph{Impact:} Gates Fermion Masses to A+.
        \item \emph{Approach:} Exploit the spectral structure of the Dirac operator on the internal manifold (likely $\mathbb{CP}^2$ or a Bieberbach manifold).
        \item \emph{Estimated timeline:} 1--6 months if the spectral structure cooperates; may require new mathematical tools.
    \end{itemize}
\end{enumerate}


%% ============================================================
\section{Comparison: Current State vs.\ All-A+ DFD}
%% ============================================================

\begin{longtable}{p{3.5cm}p{5cm}p{5cm}}
\toprule
\textbf{Sector} & \textbf{Current State} & \textbf{At A+} \\
\midrule
\endhead
Parameter Economy &
Zero free parameters, all quantities derived from $S[\psi]$. &
Same (already A+). With a complete census table and published parameter count. \\
\midrule
Strong CP &
$\bar{\theta}=0$ proven structurally. Not yet published. &
Refereed publication. Neutron EDM prediction computed and compared with data. Axion alternative explicitly contrasted. \\
\midrule
Testability &
Many predictions made; none independently tested. Three YELLOW predictions experiment-limited. &
At least one confirmed prediction. Pre-registered forecasts with DOIs. Peer-reviewed prediction papers. Active experimental collaborations. \\
\midrule
Gravitational &
Newtonian limit and Schwarzschild solution. No global existence. Horizon BCs incomplete. &
Global existence (at least symmetric case). 2PN corrections matching binary pulsars. GW templates consistent with LIGO. Rigorous horizon BCs. \\
\midrule
Galactic / MOND &
Flat rotation curves and Tully--Fisher derived. Bullet Cluster 2--3$\times$ deficit. &
175 SPARC galaxies fitted to $<$10\% residuals. Bullet Cluster resolved. RAR derived analytically. Cluster-scale lensing predictions. \\
\midrule
Microsector / $\alpha$ &
$\alpha = 1/137.036\ldots$ to 9+ digits. Step 9 and $k_{\max}$ gaps. &
Every step proven from the action. $k_{\max}$ derived from stability. External independent verification. Published in PRL/PRD. \\
\midrule
Math Foundations &
Well-posed action, local well-posedness, spectral geometry on $S^3$. Global existence open. &
Global existence (spherical + small data). Complete branch classification. Rigorous horizon BCs. Published mathematical proofs. \\
\midrule
Quantum Sector &
Decoherence mechanism described. Page--Geilker resolution exists but unformalised. Born rule not derived. &
Measurement theory derived from $S[\psi]$. Decoherence rates computed and tested. Born rule derived. Published comparison with all interpretations. \\
\midrule
Fermion Masses &
12 masses derived with $n_f$ exponents partially justified. 3.3\% electron error. &
All 12 masses to $<$1\%. $n_f$ exponents rigorously derived. Matching conditions specified. Neutrino mass ordering predicted. \\
\midrule
Cosmology &
First CMB spectrum computed (Bolt.jl). No production pipeline. $H_0$ not quantitatively computed. &
Full Boltzmann pipeline validated against Planck. $H_0$ tension resolved quantitatively. $P(k)$ and $S_8$ matched. Zero-parameter MCMC goodness-of-fit published. \\
\midrule
Comp.\ Maturity &
Working research code in Julia. No CI, no docs, no public packages. &
Documented, tested, packaged codes with DOIs. Production Boltzmann solver. External users running DFD computations independently. \\
\midrule
Unification Reach &
$\alpha$, Strong CP, fermion masses, gravity, dark matter/energy derived. CKM/PMNS incomplete. $\Lambda_{\text{QCD}}$ not derived. &
All 19 SM parameters derived. CKM and PMNS matrices predicted. $\Lambda_{\text{QCD}}$ and $\alpha_s(M_Z)$ computed. Complete unified table published. \\
\bottomrule
\end{longtable}

\noindent\textbf{Summary:} Current DFD is a theory of extraordinary breadth that derives actual numbers but lacks computational closure, formal mathematical completeness, and external validation. At all-A+, DFD would be the most predictive and parameter-free theory of fundamental physics ever constructed, with every claim independently verified and experimentally tested. The gap between these two states is approximately 2--5 years of focused work, with some items potentially requiring longer.


%% ============================================================
\section{Honest Assessment: Which Sectors May Never Reach A+?}
%% ============================================================

Scientific honesty requires acknowledging that some goals may be unachievable.

\subsection{Likely to Reach A+ (with sufficient effort)}
\begin{itemize}
    \item \textbf{Parameter Economy} --- Already there.
    \item \textbf{Strong CP} --- Gap is publication only. Extremely likely.
    \item \textbf{Computational Maturity} --- Engineering problem. Achievable with person-hours.
    \item \textbf{Microsector / $\alpha$} --- Gaps are well-defined mathematical problems with strong numerical evidence supporting the proofs.
\end{itemize}

\subsection{Achievable but Challenging}
\begin{itemize}
    \item \textbf{Gravitational / Classical} --- Post-Newtonian and GW templates are achievable (A). Global existence is hard but not necessarily impossible for the specific $\psi$-field equation (A+).
    \item \textbf{Galactic Dynamics / MOND} --- SPARC fits and RAR are achievable (A). Bullet Cluster is the wildcard.
    \item \textbf{Fermion Masses} --- The 3.3\% error is likely fixable. The $n_f$ derivation is the hard part but the spectral framework is promising.
    \item \textbf{Cosmology} --- The pipeline is buildable. The question is whether DFD actually matches Planck data when fully computed. If it does: A+. If not: the theory has a problem deeper than computational maturity.
    \item \textbf{Testability} --- Fundamentally externally gated. DFD can do everything right and still wait years for experiments.
\end{itemize}

\subsection{May Never Reach A+}
\begin{itemize}
    \item \textbf{Math Foundations.} Global existence in 3+1 dimensions is a genuinely hard open problem. The Yang--Mills mass gap (a simpler-looking problem) has been open for decades and carries a \$1M Millennium Prize. If the $\psi$-field equation has similar structural obstacles, the proof may take decades or may not exist. \textbf{Realistic ceiling: A} (with spherically symmetric global existence and complete branch analysis).

    \item \textbf{Quantum Sector.} The Born rule derivation is the sticking point. No approach to quantum mechanics has fully derived the Born rule from dynamics without some circularity. If DFD achieves this, it would be one of the most important results in the foundations of quantum mechanics. But it may not be achievable from $\psi$-field dynamics alone. \textbf{Realistic ceiling: A$-$} (with formalized Page--Geilker, decoherence rates, and clear comparison with alternatives).

    \item \textbf{Unification Reach.} Deriving the CKM and PMNS matrices from first principles has never been done by any theory. It is possible that the $\psi$-field framework does not contain enough structure to determine the 4 CKM parameters and 6 PMNS parameters uniquely. If mixing requires additional input beyond $S[\psi]$, this sector cannot reach A+ without modifying the theory. \textbf{Realistic ceiling: A$-$} (with $\Lambda_{\text{QCD}}$ derived and partial CKM results).
\end{itemize}

\subsection{The Cosmology Existential Risk}
There is a scenario that must be stated plainly: \textbf{when the full Boltzmann pipeline is built and run, DFD's cosmological predictions may not match Planck data.} If the CMB power spectrum deviates from observations by more than a few percent, this is not fixable by tweaking---DFD has no free parameters to adjust. Such a failure would not just prevent Cosmology from reaching A+; it would undermine the entire framework.

This is simultaneously DFD's greatest risk and its greatest virtue: a zero-parameter theory is maximally falsifiable. Building the Boltzmann pipeline is therefore not just a path to A+; it is an existential test of the theory. This test \emph{must} be run, regardless of the outcome.


%% ============================================================
\section{Master Timeline}
%% ============================================================

\begin{center}
\begin{tabular}{llp{7cm}}
\toprule
\textbf{Period} & \textbf{Target} & \textbf{Key Deliverables} \\
\midrule
Months 1--2 & Quick wins & Parameter census, Zenodo deposits, pre-registration, Page--Geilker draft, $\alpha$ notebook, code packaging \\
\midrule
Months 2--4 & Paper submissions & Strong CP paper submitted, $\alpha$ paper submitted, SPARC fits computed, electron error diagnosed \\
\midrule
Months 3--6 & Pipeline foundation & \texttt{mochi\_class} background solver working, spherically symmetric global existence proved, $n_f$ exponent investigation, $k_{\max}$ derivation \\
\midrule
Months 6--9 & Pipeline first results & CMB $C_\ell$ computed, $H_0$ value extracted, Step 9 proof completed, decoherence rates computed \\
\midrule
Months 9--12 & Validation & Full CMB comparison with Planck, matter power spectrum, $\Lambda_{\text{QCD}}$ derivation, CKM investigation begun \\
\midrule
Year 2 & Maturity & Pipeline validated and public, MCMC analysis published, GW templates computed, CKM/PMNS progress, external engagement \\
\midrule
Year 3+ & A+ push & Experimental results coming in, peer-reviewed publications accumulating, community adoption beginning, hard math problems progressing \\
\bottomrule
\end{tabular}
\end{center}


%% ============================================================
\section{Resource Requirements}
%% ============================================================

Reaching all-A+ (or as close as possible) requires:

\begin{enumerate}
    \item \textbf{Boltzmann pipeline developer}: 1 FTE for 6--12 months. This is the single most important hire/collaborator. Ideally someone with CLASS or CAMB experience.
    \item \textbf{PDE/mathematical analyst}: For the global existence program. Could be a collaborator rather than full-time.
    \item \textbf{Spectral geometer}: For $n_f$ exponents, CKM/PMNS, and the $B/A$ ratio proof. Could be the same person as (2) or a specialist in non-commutative geometry.
    \item \textbf{Paper writing and submission}: 3--6 months of focused writing time across multiple papers.
    \item \textbf{Computational infrastructure}: GPU cluster access for Boltzmann runs and MCMC chains. Cloud credits or university allocation.
\end{enumerate}

\noindent Without dedicated personnel, the timeline extends by a factor of 2--3$\times$. With a small team (3--4 people), the achievable A+ sectors could be closed within 18--24 months.


%% ============================================================
\section{Conclusion}
%% ============================================================

DFD is a theory of remarkable ambition and surprising success. It derives actual numbers---$\alpha$ to 9 digits, fermion masses to a few percent, galactic rotation curves without dark matter---from a zero-parameter action. No competitor matches this breadth-plus-specificity combination.

The path to A+ in every sector is partially obstructed by genuinely hard problems (global existence, Born rule, CKM/PMNS) that may take years or may prove impossible. But the majority of the gap is \emph{not} fundamental: it is engineering (Boltzmann pipeline), formalization (Page--Geilker paper, $\alpha$ Step 9), and external validation (peer review, experiments).

The recommended strategy is:
\begin{enumerate}
    \item \textbf{Immediately:} Execute all quick wins (2--3 months). These are high-impact, low-effort actions that improve multiple sectors.
    \item \textbf{Next:} Build the Boltzmann pipeline (3--8 months). This is the highest-leverage single investment.
    \item \textbf{In parallel:} Submit papers to peer review. Engage experimentalists. Begin the hard math problems.
    \item \textbf{Ongoing:} Maintain intellectual honesty. If the Boltzmann pipeline shows DFD's cosmological predictions fail, report that finding. A theory that can be falsified and \emph{is} falsified has still contributed to physics.
\end{enumerate}

The honest best case: 9 of 12 sectors at A+ within 2--3 years, with Math Foundations, Quantum Sector, and Unification Reach at A or A$-$. The realistic ceiling for the theory is not all-A+; it is ``mostly A+, with a few sectors limited by open problems that are hard for \emph{any} theory.'' That would still make DFD the most predictive unified framework in theoretical physics.

\end{document}
